\section{Introduction}\label{sec:introduction}

Optimization-based bound tightening (OBBT) improves the bounds of a variable
by minimizing and maximizing it over a convex relaxation of the feasible
set, usually together with an objective cutoff given by the incumbent value.
Tighter bounds strengthen every relaxation that depends on them, such as the
McCormick envelopes of products, and can therefore shrink a branch-and-bound
tree. A round costs up to $2n$ auxiliary optimization problems for $n$
variables, and after a round the relaxation can be rebuilt on the smaller box
and the round repeated. Solvers ration this work: they run OBBT mainly at the
root, skip directions that a known relaxation point already shows cannot
improve, order the problems to reuse warm starts, impose work limits, and
derive reusable bound inequalities from the dual
solutions~\cite{gleixner2017-three-enhancements-for-optimization-based,scip2026obbt}.
Learned variants choose which variables to tighten in each
round~\cite{cengil2025learning} or which domain-reduction configuration to
use on an instance~\cite{gomezcasares2025domain}.

Whenever OBBT is available, three decisions remain: whether to start it at a
node, when to stop repeating it, and when a change of the search state, such
as a better incumbent or a smaller node box, justifies running it again. Each
decision needs some knowledge of what further tightening can still achieve.
This paper studies that question for a specified relaxation family. We ask
what limits further OBBT, how the limit and the speed at which it is
approached depend on the cutoff and on the relaxation, and which of these
limits can be certified by finite computations.

The starting point is the order structure of the operator $T_U$ that maps a
box to the smallest box containing all of its points whose relaxed objective
is at most the cutoff $U$. For a valid relaxation family that is monotone
under box restriction, there are three limits on further tightening
(Section~\ref{sec:foundations}):
\begin{enumerate}[label=(\roman*)]
\item the box hull of the original feasible points with objective at most
  $U$, which no valid tightening removes and which shrinks with a better
  cutoff when the problem domain is held fixed;
\item the fixed boxes of $T_U$, which no number of further rounds of the same
  relaxation family removes and only a smaller cutoff, a stronger
  relaxation, or branching can shrink;
\item the rate at which the iterates approach their limit, which determines
  how many rounds are needed to reach a given movement tolerance.
\end{enumerate}
The first is a property of the problem, the second of the relaxation, and the
third of the iteration. The paper develops each of them.

\paragraph{Local rates and stalling.} Near a minimizer $x^*$ we rescale the
relaxation to boxes $x^*+tD(d)$ of a fixed shape $d$ and let $t\to0$
(Section~\ref{sec:local-rates}). The limit is a second-order tangent model
$Q(d,\xi)$ that depends on the shape of the box and on the position $\xi$ in
it. It defines an order-preserving, positively homogeneous map $\Phi$ on box
shapes whose upper Collatz--Wielandt factor $r^*$ bounds the local linear
rate: if $r^*<1$, then for every $\lambda\in(r^*,1)$ the iterates from a
sufficiently small box around $x^*$ contract by the factor $\lambda$ per
round, in a suitable gauge, until their width is $O(\sqrt\epsilon)$, where
$\epsilon=U-f^*$ is the cutoff slack, and the relaxation gap that remains is
$O(\epsilon)$. If points with negative tangent value attain all faces of a
shape, boxes of that shape around $x^*$ are fixed at every sufficiently
small scale and every cutoff $U\ge f^*$, so OBBT stalls there. For
quadratic objectives with McCormick relaxations the tangent model is exact.
This gives closed forms: the rate $(\sqrt{2a^2+4\abs a}-\abs a)/2$ for
$x^2+y^2+axy$ with $0<\abs a<2$, a continuum of eigenvector boxes, the exact
limit and remaining gap at a positive cutoff slack, row conditions for
stalling and for contraction, and a strongly convex family that stalls
exactly when $a(n-1)(n-2)\ge2$. For
composite McCormick relaxations of factorable functions with twice
continuously differentiable univariate factors, we prove a second-order
expansion that computes $Q$ by a recursion over the factors. At minimizers
on the boundary of the box with strict complementarity, the free coordinates
obey the bounds of the interior theory, and the active widths of the limit
are $O(\epsilon)$.
These conditions are sufficient, not exhaustive: we give a valid monotone
family that converges sublinearly and whose limit has width of order
$\epsilon^{1/3}$.

\paragraph{Finite certificates.} A finite set of relaxation points whose box
hull is a box $P$, and which are feasible in the relaxation rebuilt on $P$,
proves that $P$ is fixed (Section~\ref{sec:certificates}). It therefore
bounds the total movement of every endpoint under all later rounds and
directional updates, at every cutoff not below the largest relaxation value
of the points, as long as the box being tightened contains $P$. Established feasibility filtering, which
checks the points only in the current relaxation, bounds a single frozen
round. We give exact cutoff thresholds and the adaptation of cached witness
pools to a new cutoff (Section~\ref{sec:cutoff}), objective ceilings, the
scope of the certificates under integer rounding, cuts, and branching, and
the difference between checking a certificate and finding one. When the
iterates converge only asymptotically, the residual of one exact round and a
verified uniform sensitivity matrix bound the remaining movement
(Section~\ref{sec:residual}).

\paragraph{Constraints.} Fixed-matrix LP duality and exact basis regions
bound how the coordinate supports respond to a cutoff change across
active-set changes; a complete cover by basis regions supplies the uniform
sensitivity matrix that the residual certificates need; and a
feasible-repair argument transfers relaxation and constraint errors into a
contraction bound for nonlinear constraints, with explicit constants in a
worked example (Section~\ref{sec:constraints}).

\paragraph{Algorithms, accounting, and evidence.} An exact reference procedure
commits a round only after verifying a primal--dual optimality proof for
every endpoint, and reports a fixed box only when the finite certificate
holds. A numerical policy
inside SCIP applies dual bounds evaluated with directed rounding, screens
directions with points of the current relaxation, stops unproductive pilots,
and reconsiders tightening after the incumbent improves
(Section~\ref{sec:algorithms}). A cost ledger bounds optional work on the
enhanced execution but cannot bound its run time relative to an unchanged
baseline (Section~\ref{sec:effort}). Section~\ref{sec:experiments} reports
the computational evidence.

The rate theory and the certificates answer different questions and
complement each other. The rate theory identifies which limit dominates near
a minimizer and how it scales with the cutoff, but its hypotheses involve
$x^*$, $f^*$, and tangent data that a solver does not know. The certificates
are computable once their witnesses or sensitivity data are available. A
certificate for the current box proves exact fixedness; a certificate for
an inner box bounds the remaining movement and can justify stopping at a
specified tolerance. Neither requires positive width: a singleton is also
a fixed box if its point belongs to its own cutoff relaxation. The rate
theory describes which fixed boxes can occur locally. In the stall regime,
positive-width fixed boxes exist at every sufficiently small scale and
finitely many witnesses certify them. In the strictly contracting regime,
under the closedness property of Section~\ref{sec:foundations}, the limit is
a fixed box of width $O(\sqrt\epsilon)$ and at most $2n$ witnesses certify
it. At $\epsilon=0$ the only nearby fixed box is $\{x^*\}$. If an iterate
reaches it, a finite certificate proves termination. If convergence is only
asymptotic, the singleton of an optimal incumbent in the current box is
still protected (Corollary~\ref{cor:original-points}); it bounds the
remaining movement, and the bound is exact if the iterates converge to that point
(Example~\ref{ex:sharp-halving}). Otherwise a bound on the remaining
movement, such as a residual certificate, justifies a tolerance stop. The
upper bounds of order $\sqrt\epsilon$ on the width of the limit and of
order $\epsilon$ on the remaining gap motivate reconsidering OBBT after the
incumbent improves; in the two-variable model the limit itself shrinks at
exactly these orders. A stored certificate also
states exactly when its own proof expires: when the cutoff falls below the
face threshold of its pool. The certified box can remain fixed below that
threshold, down to an exact box threshold computed by at most $2n$ face
optimizations (Section~\ref{sec:cutoff}).

The computational evidence limits what these guarantees imply. In a
prospective comparison of 120 runs on twenty models with two seeds each, an
adaptive policy used $20.9\%$ fewer auxiliary LPs than a fixed-order policy
but spent more time in its propagator, and neither added policy produced an
additional solve or a demonstrated speedup over native SCIP within ten
seconds. An earlier study of iterated OBBT as an external presolve on 339
MINLPLib instances found that the tightened boxes reduced node counts, but
the savings in the final solve did not offset the cost of OBBT. The
measured policy used only current-round screening, not the all-future
certificates, so the experiments say nothing about the computational value
of those certificates. They show that the number of auxiliary LPs is not a
cost measure, and Section~\ref{sec:effort} shows why a bound on width or on
the remaining tightening cannot by itself bound search time. The certificates
give finite stopping guarantees for a specified tightening operator; whether
acting on them saves time is a separate question that requires evidence of
its own.

\paragraph{Attribution.} The order properties of the operator, the
persistence of fixed boxes, and the existence of a greatest fixed box are
elementary consequences of monotonicity; the last is an instance of
Tarski's fixed-point theorem~\cite{tarski1955fixpoint}. The limits of
iterated domain reduction and their order independence were studied by
Caprara and
Locatelli~\cite{caprara2010-global-optimization-problems-and-domain}, and
Caprara, Locatelli, and Monaci gave examples in which iterated reduction
does not move any
bound~\cite{caprara2016-theoretical-and-computational-results-about}. The
scalar contraction bounds follow the mechanism of marginals-based range
reduction~\cite{ryoo1995-global-optimization-of-nonconvex-nlps}, and their
strict contraction condition has the same numerical threshold as known
nonstrict conditions that exclude clustering in branch and
bound~\cite{wechsung2014-the-cluster-problem-revisited,kannan2017-the-cluster-problem-in-constrained}.
Repeated OBBT is also used in adaptive partitioning and power-flow
optimization~\cite{nagarajan2019-an-adaptive-multivariate-partitioning-algorithm,sundar2023-optimization-based-bound-tightening-using}.
Convergence-order analysis of McCormick and McCormick--Taylor relaxations is
established~\cite{bompadre2012-convergence-rate-of-mccormick-relaxations,najman2016-convergence-analysis-of-multivariate-mccormick-relaxations,bompadre2013-convergence-analysis-of-taylor-models}.
Feasibility filtering, dual reuse, and incumbent-triggered reevaluation of
stored bounds exist in prior work and in
SCIP~\cite{gleixner2017-three-enhancements-for-optimization-based,scip2026obbt}.
LP duality, parametric LP, and contraction principles are classical. Our
contribution consists of the OBBT-specific formulations: the tangent model
and tangent map with their contraction and stall theorems, the
relaxation-specific exact rates and stall conditions, the explicit
shape-dependent second-order expansion of composite McCormick relaxations,
the finite certificate formulations together with their scope, and a
prospective evaluation of an adaptive scheduling policy that uses only
current-round screening and heuristic triggers, not the all-future
certificates. The local results concern the
shape and position dependent tangent map of a specified relaxation family;
they do not assert priority for iterated domain reduction, adaptive OBBT,
or convergence-order analysis in general.

\paragraph{Organization.} Section~\ref{sec:related} reviews related work.
Section~\ref{sec:foundations} defines the operator and its limits, and
Section~\ref{sec:local-rates} develops the local theory.
Sections~\ref{sec:certificates}, \ref{sec:cutoff}, and~\ref{sec:residual}
develop the certificates, and Section~\ref{sec:constraints} treats
constraints. Sections~\ref{sec:algorithms} to~\ref{sec:experiments} describe
the implementations, the cost accounting, and the computational evidence,
and Section~\ref{sec:discussion} discusses the consequences for starting,
stopping, and repeating OBBT. Appendices~\ref{sec:local-proofs}
and~\ref{sec:parametric-proofs} contain longer proofs,
Appendix~\ref{sec:numerical-validation} gives the numerical bound validation,
Appendix~\ref{sec:scheduling} gives scheduling guarantees that keep the
baseline solver unchanged, and Appendix~\ref{sec:precursor-details}
describes the precursor study.
